\chapter{Méthodologie expérimentale et évaluation}
\label{chap:methodologie}

\section{Extraction des caractéristiques}

SVM, Random Forest et k-NN ne lisent pas directement une image : il faut d'abord la résumer en une liste de nombres, le \emph{descripteur}. Nous utilisons \textbf{HOG} (\emph{Histogram of Oriented Gradients}), qui décrit l'orientation des contours :

\begin{enumerate}
  \item l'image 64~×~64 est convertie en niveaux de gris ;
  \item elle est découpée en cellules de 8~×~8~px, et chaque cellule compte ses contours selon 9 orientations ;
  \item les cellules sont regroupées en blocs de 2~×~2, normalisés pour compenser les différences d'éclairage.
\end{enumerate}

Un bloc de 16~px, décalé de 8~px à chaque fois, prend 7 positions par côté, ce qui donne $7 \times 7 = 49$ blocs de $4 \times 9$ valeurs : $49 \times 4 \times 9 = 1\,764$ nombres par image. Nous testerons aussi des cellules de 4~px, plus fines mais beaucoup plus nombreuses. HOG ignore la couleur : c'est une limite connue, qui pèsera surtout sur des classes de formes proches, comme le bus et le camion.

\section{Présentation des algorithmes SVM, Random Forest et k-NN}

\begin{center}
\small
\begin{tabularx}{\textwidth}{>{\raggedright\arraybackslash}p{2.3cm}YY>{\raggedright\arraybackslash}p{3cm}}
  \toprule
  \textbf{Algorithme} & \textbf{Principe} & \textbf{Points forts et limites} & \textbf{Réglages testés} \\
  \midrule
  SVM & Cherche la frontière qui sépare les classes avec la plus grande marge. Avec un noyau RBF, cette frontière peut être courbe & Efficace en grande dimension ; lent à entraîner quand les exemples sont nombreux & Linéaire ou RBF ; paramètre C \\
  Random Forest & Fait voter des centaines d'arbres de décision, chacun entraîné sur un tirage différent des données & Rapide et robuste ; chaque question d'un arbre ne regarde qu'une valeur à la fois & Nombre d'arbres, profondeur \\
  k-NN & Donne la classe majoritaire parmi les $k$ images d'entraînement les plus proches & Simple, sans vrai entraînement ; les distances perdent leur sens en grande dimension & $k$ ; réduction par PCA (50 à 200 dimensions) \\
  \bottomrule
\end{tabularx}
\end{center}

Le SVM et le k-NN sont précédés d'une mise à l'échelle des valeurs, car ils y sont sensibles. Le Random Forest n'en a pas besoin.

\textbf{Extensions, sous réserve de l'accord de l'encadrant.} En plus des trois algorithmes du sujet, nous proposons XGBoost, une autre famille d'arbres construits l'un après l'autre, et un petit réseau convolutif (CNN) entraîné de zéro. Ces deux extensions sont comparées de la même façon, sans remplacer les trois modèles demandés.

\section{Protocole expérimental}

\begin{enumerate}
  \item \textbf{Même point de départ :} tous les modèles utilisent les mêmes recadrages, le même découpage et les mêmes descripteurs.
  \item \textbf{Réglages sur la validation :} chaque combinaison de réglages est entraînée sur le train, puis mesurée sur la validation.
  \item \textbf{Une seule chose change à la fois :} chaque essai est noté avec ses réglages, son score et son temps d'entraînement.
  \item \textbf{Test ouvert une seule fois :} les modèles retenus sont mesurés sur le test à la fin. Plus aucun réglage n'est permis après.
  \item \textbf{Reproductibilité :} les graines aléatoires sont fixées, et les notebooks s'exécutent de bout en bout.
\end{enumerate}

Le test reste fermé jusqu'à la fin, car un modèle réglé en regardant le test aurait un score trop optimiste. C'est comme un élève qui connaîtrait les questions de l'examen.

\section{Métriques d'évaluation}
\label{sec:metriques}

\begin{center}
\small
\begin{tabularx}{\textwidth}{>{\raggedright\arraybackslash}p{3.2cm}Y}
  \toprule
  \textbf{Métrique} & \textbf{Ce qu'elle mesure} \\
  \midrule
  Accuracy & Part de bonnes réponses sur toutes les images \\
  Précision (par classe) & Quand le modèle répond « camion », part de vrais camions \\
  Rappel (par classe) & Parmi les vrais camions, part que le modèle retrouve \\
  F1 (par classe) & Moyenne harmonique de la précision et du rappel : $F1 = \frac{2 \times \text{précision} \times \text{rappel}}{\text{précision} + \text{rappel}}$ \\
  F1 macro & Moyenne des F1 des 6 classes, chacune avec le même poids : \textbf{métrique principale} \\
  Rappel « personne » & Part des personnes reconnues : l'erreur la plus grave pour un véhicule \\
  Matrice de confusion & Croise vraies classes et classes prédites, pour voir quelles classes se confondent \\
  Temps d'inférence & Temps de calcul d'une prédiction sur CPU \\
  \bottomrule
\end{tabularx}
\end{center}

\section{Méthode de comparaison des modèles}

\begin{itemize}
  \item \textbf{Mêmes données :} les modèles sont comparés sur les mêmes images, dans le même ordre.
  \item \textbf{Scores planchers :} un modèle qui répond toujours la même classe, ou qui tire au hasard selon les fréquences, donne la valeur à dépasser.
  \item \textbf{Analyse des erreurs :} au-delà du score global, nous comparons les matrices de confusion et regardons des images mal classées tirées au hasard.
  \item \textbf{Incertitude :} avec environ 2~250 images de test, une accuracy n'est connue qu'à environ $\pm$~2 points près. Un écart plus petit ne permet pas de départager deux modèles.
\end{itemize}

\section{Critères de sélection du modèle final}

\begin{enumerate}
  \item le meilleur \textbf{F1 macro sur la validation} ;
  \item à score proche, le meilleur \textbf{rappel « personne »} ;
  \item puis un \textbf{temps d'inférence} compatible avec le simulateur (quelques millisecondes) ;
  \item enfin, le modèle le plus simple et le plus stable d'un essai à l'autre.
\end{enumerate}

Le modèle retenu est exporté avec ses métadonnées. Son seuil « Inconnu » est choisi sur la validation : c'est le plus haut seuil qui laisse au moins 90~\% des images classées.

\section{Protocole de validation dans le simulateur}

Chaque classifieur parcourt exactement les mêmes épisodes : même carte, même itinéraire, mêmes obstacles et mêmes images. Le protocole prévoit 50 épisodes, avec une graine fixée.

\begin{center}
\small
\begin{tabularx}{\textwidth}{>{\raggedright\arraybackslash}p{4.2cm}Y}
  \toprule
  \textbf{Mesure} & \textbf{Définition} \\
  \midrule
  Épisodes arrivés & Le véhicule a atteint la destination \\
  Collisions & Le véhicule a atteint un obstacle encore présent \\
  Obstacles bien classés & Classe prédite égale à la vraie classe \\
  Fautes graves & Personne classée autrement que « personne », « Inconnu » compris \\
  Décisions inadaptées & Action différente de celle que demandait la vraie classe \\
  Détour et recalculs & Distance parcourue rapportée au trajet prévu, nombre de nouveaux calculs \\
  Temps par décision & Temps moyen d'une prédiction, mesuré sur CPU \\
  \bottomrule
\end{tabularx}
\end{center}

Le mode « simulation » n'est pas un modèle : son taux d'erreur est réglé à la main. Il vérifie seulement que la logique du simulateur fonctionne, et ses chiffres ne se comparent pas à ceux d'un vrai classifieur.
